\documentclass[11pt,twocolumn]{article}
\usepackage[margin=0.85in,columnsep=0.25in]{geometry}
\usepackage{booktabs}
\usepackage{array}
\usepackage{hyperref}
\usepackage{parskip}
\usepackage{multirow}

\title{\textbf{581 Report --- Auditing LLM Summaries:\\Sprint 3 Multi-Task Learning}}

\author{Jai Sharma \and Gilly Guo \and Rayan Jaz \and Sina Oloomi\\[0.5em]
\small University of British Columbia}

\date{}

\begin{document}

\maketitle

% ─────────────────────────────────────────────────────────────────────────────
\section{Abstract}
% ─────────────────────────────────────────────────────────────────────────────

Sprint 3 extends the Sprint 2 transfer learning pipeline by introducing a
multi-task learning (MTL) objective. Instead of predicting only the
\texttt{accuracy} label, models now jointly predict \texttt{accuracy} and
\texttt{brevity} as a single combined label. We evaluate a traditional
(Logistic Regression) and neural (MLP) MTL model, then compare several
ensemble rules that combine Sprint 2 and Sprint 3 predictions. The Sprint 3
traditional MTL achieves the best pure Macro F1 of 0.421, outperforming the
Sprint 2 weighted ensemble (0.372), primarily by recovering recall on the
minority class 1. A confidence-gated hybrid of Sprint 2 and Sprint 3
predictions offers the best overall balance, reaching 0.750 accuracy and 0.415
Macro F1.

% ─────────────────────────────────────────────────────────────────────────────
\section{Introduction}
% ─────────────────────────────────────────────────────────────────────────────

Sprint 1 established TF-IDF baselines and Sprint 2 introduced GloVe transfer
learning and soft-voting ensembles. A persistent challenge across both sprints
was severe class imbalance, which caused models to over-predict class 2 with
near-zero recall on classes 0 and 1. Sprint 3 addresses this by leveraging the
\texttt{brevity} annotation as an auxiliary signal, with the hypothesis that
brevity and accuracy are correlated and that jointly modeling both may improve
minority class recovery.

Multi-task learning is a natural fit for this problem. Both labels are human
annotations over the same summaries, so the model can share representations
while simultaneously learning to distinguish multiple aspects of summary
quality. Rather than modifying the input features as in Sprint 2, Sprint 3
changes the prediction objective itself, allowing the auxiliary task to act as
a regularizer on the primary task.

% ─────────────────────────────────────────────────────────────────────────────
\section{Data}
% ─────────────────────────────────────────────────────────────────────────────

The same 400-entry CMV corpus and 320/40/40 train/dev/test split from Sprint 1
is used in this sprint, partitioned by thread and stratified by reviewer to
prevent data leakage. No new data was collected. The key change is that both
the \texttt{accuracy} (0--2) and \texttt{brevity} (0--2) annotation labels are
now used during training, with \texttt{accuracy} remaining the primary
evaluation target and \texttt{brevity} serving as the auxiliary task.

\subsection{Auxiliary Task Selection}

\texttt{brevity} was chosen as the auxiliary task over \texttt{sentiment} for
two reasons. First, \texttt{brevity} has a healthier class spread than
\texttt{sentiment}, which is heavily imbalanced, producing a more complete
joint label space. Second, in a direct comparison, \texttt{accuracy} +
\texttt{brevity} gave a better Macro F1 than \texttt{accuracy} +
\texttt{sentiment} on the traditional multi-task setup. The intuition is that a
summary capturing most key points concisely is often also more accurate, meaning
brevity provides a complementary signal about overall summary quality that is
closely related to the primary task.

\subsection{Joint Label Construction}

Labels from both tasks are combined into a single string of the form
\texttt{acc\_\{accuracy\}\_\_brev\_\{brevity\}} (e.g.\ \texttt{acc\_2\_\_brev\_1}),
producing a set of joint classes that the model predicts as a single multi-class
problem. After prediction, the joint label is decoded back into separate
\texttt{accuracy} and \texttt{brevity} predictions for evaluation. A
\texttt{LabelEncoder} is fit on the training joint labels and applied
consistently to the dev set.

% ─────────────────────────────────────────────────────────────────────────────
\section{Methodology}
% ─────────────────────────────────────────────────────────────────────────────

\subsection{Feature Representation}

Feature construction is identical to Sprint 2. The source text and summary are
concatenated as \texttt{SOURCE: ...\textbackslash nSUMMARY: ...} to form a
unified input. TF-IDF features are extracted using unigrams and bigrams, with
English stop words removed, a minimum document frequency of 2, and a vocabulary
capped at the top 5,000 features. GloVe embeddings
(\texttt{glove-wiki-gigaword-50}) are used to generate separate mean embeddings
for the source and summary, which are then concatenated to form a 100-dimensional
dense representation.

\subsection{Traditional MTL Model}

The traditional MTL model horizontally stacks the sparse TF-IDF features and
the dense GloVe features, then trains a Logistic Regression classifier on the
joint labels with \texttt{max\_iter=2000}, \texttt{C=2.0}, and balanced class
weights. After inference, the joint predictions are decoded and the
\texttt{accuracy} component is extracted for evaluation against the gold labels.

\subsection{Neural MTL Model}

The neural MTL model reduces the TF-IDF features to 50 components via
Truncated SVD, down from 100 in Sprint 2, and concatenates the result with the
GloVe features before passing them to a single hidden layer MLP with 64 neurons,
\texttt{max\_iter=500}, and early stopping enabled. The same joint label
training and decoding procedure is used as in the traditional model.

\subsection{Ensemble Experiments}

Six combination rules are evaluated using saved dev predictions from the Sprint
2 weighted ensemble and the Sprint 3 traditional MTL model. The Sprint 2 and
Sprint 3 predictions are used as individual baselines. The min-label rule takes
the lower of the two predictions. The brevity-gate rule uses the Sprint 3
prediction only when predicted brevity is low ($\leq$1) and falls back to
Sprint 2 otherwise. The confidence-gate rule uses Sprint 3 when its joint-label
probability margin exceeds 0.329 and falls back to Sprint 2 otherwise. Finally,
a three-model majority vote adds the Sprint 1 traditional baseline as a third
voter, with ties broken by taking the lower label. The best rule is selected by
Macro F1 with accuracy as a tiebreaker.

% ─────────────────────────────────────────────────────────────────────────────
\section{Results}
% ─────────────────────────────────────────────────────────────────────────────

\subsection{Full Model Comparison Across Sprints}

\begin{table}[ht]
\centering
\small
\begin{tabular}{lcc}
\toprule
\textbf{Model} & \textbf{Acc.} & \textbf{F1} \\
\midrule
S1 traditional baseline   & 0.650 & 0.336 \\
S2 traditional transfer   & 0.625 & 0.366 \\
S2 weighted ensemble      & 0.725 & 0.372 \\
Neural MTL (MLP)          & 0.700 & 0.275 \\
Traditional MTL (LogReg)  & 0.700 & \textbf{0.421} \\
Confidence-gated (S2+S3)  & \textbf{0.750} & 0.415 \\
\bottomrule
\end{tabular}
\caption{Dev performance across all sprints.}
\end{table}

\subsection{Ensemble Rule Breakdown}

\begin{table}[ht]
\centering
\small
\begin{tabular}{lcc}
\toprule
\textbf{Rule} & \textbf{Acc.} & \textbf{F1} \\
\midrule
Sprint 2 only             & 0.725 & 0.372 \\
Sprint 3 only             & 0.700 & 0.421 \\
Min-label                 & 0.700 & 0.421 \\
Brevity-gate              & 0.700 & 0.384 \\
Confidence-gate           & 0.725 & 0.372 \\
Majority vote (3 models)  & 0.675 & 0.348 \\
\bottomrule
\end{tabular}
\caption{Ensemble rules on the Sprint 3 dev set.}
\end{table}

\subsection{Per-class Breakdown}

\begin{table}[ht]
\centering
\small
\begin{tabular}{ccccc}
\toprule
\textbf{Model} & \textbf{Cl.} & \textbf{P} & \textbf{R} & \textbf{F1} \\
\midrule
\multirow{3}{*}{Trad. MTL}
  & 0 & 0.000 & 0.000 & 0.000 \\
  & 1 & 0.455 & 0.455 & 0.455 \\
  & 2 & 0.793 & 0.821 & 0.807 \\
\midrule
\multirow{3}{*}{Neural MTL}
  & 0 & 0.000 & 0.000 & 0.000 \\
  & 1 & 0.000 & 0.000 & 0.000 \\
  & 2 & 0.700 & 1.000 & 0.824 \\
\bottomrule
\end{tabular}
\caption{Per-class precision (P), recall (R), and F1 for both MTL models.}
\end{table}

% ─────────────────────────────────────────────────────────────────────────────
\section{Analysis and Discussion}
% ─────────────────────────────────────────────────────────────────────────────

\subsection{Multi-task Learning Improves Minority Class Recall}

The traditional MTL model raises class 1 F1 from approximately 0.286 in the
Sprint 2 ensemble to 0.455, driven by the regularizing effect of the brevity
auxiliary signal. The joint label forces the model to distinguish more
fine-grained combinations of accuracy and brevity, which indirectly discourages
over-prediction of the majority class. The overall effect resembles
regularization rather than a fundamental behavioural change: the model becomes
less overconfident about class 2, which helps on borderline class 1 examples
while still correctly retaining most class 2 predictions.

\subsection{Neural MTL Collapses to Class 2}

The MLP model predicts class 2 for every example, producing a degenerate Macro
F1 of 0.275. This is consistent with the neural model's behaviour in Sprint 1
and Sprint 2. The most likely causes are the smaller feature space from reducing
SVD components to 50, higher model capacity relative to dataset size, and the
imbalanced joint label distribution overwhelming the minority signal. The
multi-task objective does not help the neural model break out of the
majority-class pattern the way it does for the linear model.

\subsection{Ensemble Results and the Confidence-Gated Hybrid}

No simple ensemble rule outperforms the Sprint 3 traditional MTL on Macro F1
alone. However, the confidence-gate rule produces the best overall balance,
reaching 0.750 accuracy and 0.415 Macro F1 by using the Sprint 3 prediction
when the model's joint-label probability margin is high and falling back to the
Sprint 2 ensemble otherwise. This is the highest accuracy result across all
three sprints and still outperforms the Sprint 2 weighted ensemble on Macro F1,
making it the most well-rounded combined model. The three-model majority vote
performs worst overall, suggesting the Sprint 1 baseline adds noise rather than
a useful third signal.

\subsection{Qualitative Error Analysis}

Comparing the Sprint 1 traditional baseline against the Sprint 3 traditional
MTL on dev predictions reveals a clear pattern in both the fixed and broken
cases \cite{claude_analysis}. Among the fixed examples, \texttt{152\_u} and
\texttt{159\_u} were cases where the baseline over-predicted class 2 and the
MTL model correctly moved predictions down to class 1. Both summaries are fairly
complete but not clearly strong enough for the top accuracy label, suggesting
the brevity signal is acting as a regularizer. A third fixed example,
\texttt{45\_s}, shows the opposite direction: the baseline under-predicted at
class 1 while the MTL model correctly moved it up to class 2, demonstrating
that the joint label can also support upward corrections on summaries that are
both accurate and complete.

Among the broken examples, \texttt{159\_s} and \texttt{184\_u} were cases where
the baseline correctly predicted class 2 but the MTL model dropped predictions
to class 1. The model's increased caution, while beneficial overall, leads to
underprediction on summaries that are genuinely strong. The brevity signal may
have pulled the joint prediction toward a lower accuracy class in these cases.

\subsection{Class 0 Remains Undetected}

No model across any sprint has successfully predicted class 0. This reflects
both the extreme rarity of class 0 examples in the training set and the
inherent difficulty of distinguishing ``not accurate'' summaries from
``moderately accurate'' ones using surface-level lexical and embedding features
alone.

\subsection{Data Limitations and Noise}

The dataset is small and the \texttt{accuracy} label is imbalanced, which
places a hard ceiling on how much any model can realistically improve. The
annotations are subjective, reflecting human judgments of summary quality rather
than objective ground truth. Furthermore, the prediction target is itself a
rating of LLM-generated summaries, adding an additional layer of noise on top
of already messy natural language data. Some instability across models and
limited absolute gains are therefore expected consequences of the low-resource
noisy setup rather than indicators of poor modeling.

% ─────────────────────────────────────────────────────────────────────────────
\section{Conclusion}
% ─────────────────────────────────────────────────────────────────────────────

Sprint 3 demonstrates that multi-task learning with \texttt{brevity} as an
auxiliary task meaningfully improves accuracy classification, particularly for
the difficult class 1. The traditional MTL model, with a Macro F1 of 0.421, is
the best single model across all three sprints. The confidence-gated hybrid, at
0.750 accuracy and 0.415 Macro F1, is the best balanced combined result and the
strongest overall model when both metrics are considered together.

Ensemble combinations of Sprint 2 and Sprint 3 models offer limited gains over
the MTL model alone, suggesting that the MTL objective is itself the more
effective integration strategy compared to post-hoc rule-based combination.
Class 0 prediction remains an open problem that will likely require more data,
oversampling, or a fundamentally different modeling approach. Future work could
explore sentence-level embeddings such as SBERT, explicit class-rebalancing via
SMOTE, or fine-tuning a small pretrained encoder directly on the annotation
task.

% ─────────────────────────────────────────────────────────────────────────────
\begin{thebibliography}{9}

\bibitem{chatgpt}
ChatGPT.
Response to ``Help brainstorm, debug, and revise project code and documentation
for this project.''
OpenAI, 18 Apr.\ 2026.
\url{https://chat.openai.com/}.

\bibitem{claude_analysis}
Claude Sonnet 4.6.
Response to ``Run the project scripts and compare the output metrics and model
behaviour against the notes I have taken. Identify if any information is 
missing between the documented results and what the scripts actually
produce.''
Anthropic, 18 Apr.\ 2026.
\url{https://claude.ai/code}.

\bibitem{claude_structure}
Claude Sonnet 4.6.
Response to ``Organise the following bullet points byfollowing the section 
headings and academic layout of the provided scientific report.''
Anthropic, 18 Apr.\ 2026.
\url{https://claude.ai/code}.

\bibitem{claude_writing}
Claude Sonnet 4.6.
Response to ``This is a draft of a scientific paper. identify where the writing
style is inappropriate for a scientific paper, and suggest corrections.''
Anthropic, 18 Apr.\ 2026.
\url{https://claude.ai/code}.

\end{thebibliography}

\end{document}
